\structural{ВСТУП}

Аеропорт із недостатньою кількістю посадкових смуг втрачає гроші двома шляхами: літаки, що довго чекають посадки, доводиться дозаправляти, а частина літаків, не дочекавшись, відлітає в інші аеропорти. Будівництво нових смуг усуває ці втрати, але коштує значних грошей, і кожна наступна смуга дає менший приріст прибутку, ніж попередня. Тому питання, скільки смуг збудувати, має кількісну відповідь: таку кількість, за якої вкладені кошти повертаються найшвидше.

Аналітично цю задачу розв'язати важко. Надходження літаків випадкові, кожен літак має власні таймери дозаправки й відльоту, що відраховуються від моменту його прибуття, а дохід і витрати залежать від того, як довго чекав саме цей літак. Класичні формули теорії масового обслуговування не враховують одночасно нетерплячих клієнтів з двома порогами очікування й вартісні показники. Імітаційне моделювання дозволяє відтворити процес функціонування системи в часі й отримати статистичні оцінки потрібних показників без спрощувальних припущень~\cite{stetsenko2010, tomashevskyi2005}.

\textbf{Тема роботи} -- імітаційна модель системи посадки літаків в аеропорту на основі формального опису стохастичною мережею Петрі.

\textbf{Мета роботи} -- визначити кількість посадкових смуг, за якої час окупності додаткових смуг $T = S/p$ мінімальний, де $S$ -- витрати на будівництво смуг, $p$ -- річний прибуток.

Для досягнення мети потрібно:
\begin{itemize}
\item описати концептуальну модель системи: процес, вхідні та вихідні змінні;
\item побудувати формальний опис моделі стохастичною мережею Петрі й визначити формули вихідних характеристик;
\item розробити й реалізувати алгоритм імітації мережі Петрі та верифікувати модель;
\item визначити тривалість перехідного періоду, тривалість прогону і кількість повторних прогонів;
\item провести експериментальне дослідження для 0--10 додаткових смуг, знайти оптимальну кількість і сформулювати пропозиції.
\end{itemize}

Роботу виконано за рекомендаціями посібника~\cite{stetsenko2024}. Формалізмом опису обрано стохастичну мережу Петрі з часовими затримками. Літаки в моделі представлено маркерами-об'єктами з ідентифікатором і моментом прибуття, що близько до ідеї кольорових мереж Петрі~\cite{jensen2009}: так кожен літак отримує власні таймери, а мережа лишається компактною.

Модель і алгоритм імітації реалізовано мовою TypeScript~\cite{typescript} у середовищі виконання Bun~\cite{bun} без сторонніх бібліотек імітаційного моделювання. Програма будує схему мережі з того самого опису, що й модель, проводить усі експерименти й будує графіки для звіту.
